\chapter{Кватернионы}
\section{Операции над кватернионами}

В рамках истории, развитие чисел шло по пути увеличения размерности или как еще упоминается в источнике ~\cite{Kuipers:1999} --- <<ранга>>:
\begin{itemize}
    \item Действительные числа имеют размерность равную единице, поле является коммутативным и есть обратные.
    \item Комплексные числа имеют размерность равную двум и само поле обладает аналогичными свойствами, что и действительные числа.
    \item Кватернионы же не являются полем, а некоммутативным телом.
    \item Еще менее известной структурой являются октавы или как их еще называют октонионы, имеющие размерность равную восьми.
\end{itemize}
Заметим, что когда мы переходим от размерности двух к размерности четырех, теряется коммутативность умножения, но для любого ненулевого числа сохраняется существование обратного элемента.Например октонионы или в целом системы, обладающие размерностью больше четырех, более не обладают ассоциативностью.

Есть основное соотношение в алгебре кватернионов:
\begin{equation*}
    i^2 = j^2 = k^2 = i j k = -\sigma,
\end{equation*}
здесь, как и ранее упоминалось, $\sigma$ --- скалярная единица.В других источниках может обозначаться как просто $-1$. Данное сооттношение, также называют правилом Гамильтона.

Сам кватернион можно задать двумя способами. Первый способ заключается в том, что мы представляем кватернион как упорядоченную четверку действительных чисел:
\begin{equation*}
    q = (q_0, q_1, q_2, q_3), \quad q_i \in \mathbb{R}.
\end{equation*}
Или можно задать как сумму скаляра и вектора:
\begin{equation*}
    q = q_0 \sigma + \vb{q} = q_0 \sigma + q_1 i + q_2 j + q_3 k.
\end{equation*}
С точки зрения линейной алгебры такая сумма будет выглядеть странно, но именно благодаря такой сумме кватернионы применимы для вращений. 

К основным операциям над кватернионами относят:
\begin{enumerate}
    \item Два кватерниона вида $p = p_0 \sigma + p_1 i + p_2 j + p_3 k$ и $q = q_0 \sigma + q_1 i + q_2 j + q_3 k$ равны только, когда равны их компоненты:
        \begin{equation*}
            p_0 = q_0, \quad p_1 = q_1, \quad p_2 = q_2, \quad p_3 = q_3.
        \end{equation*}
    \item Сложение выполняется покомпонентно:
        \begin{equation}
            p + q = (p_0 + q_0)\sigma + (p_1 + q_1)i + (p_2 + q_2)j + (p_3 + q_3)k.
            \label{eq:(1)}
        \end{equation}
    \item Умножение кватерниона на действительное число $c$:
        \begin{equation*}
            q = (c q_0)\sigma + (c q_1)i + (c q_2)j + (c q_3)k.
        \end{equation*}
    \item Норма кватерниона:
        \begin{equation}
            |q| = \sqrt{q^* q} = \sqrt{q q^*}.
            \label{eq:(2)}
        \end{equation}
    \item Обратный кватернион:
        \begin{equation}
            q^{-1} = \frac{q^*}{|q|^2}.
            \label{eq:(3)}
        \end{equation}
\end{enumerate}

Также отдельно хочу отметить сложение кватерниона, которое выполняется по формуле ~\eqref{eq:(1)}. Здесь можно выделить ряд свойств. которые аналогичны свойствам действительных чисел:
\begin{itemize}
    \item Замкнутость --- в результате суммы двух кватернионов получается снова кватернион.
    \item Нулевой элемент --- $0 = 0\sigma + 0i + 0j + 0k$.
    \item Противоположный элемент --- $-q = -q_0\sigma - q_1 i - q_2 j - q_3 k$.
    \item Ассоциативность --- $(p + q) + r = p + (q + r)$.
    \item Коммутативность --- $p + q = q + p$.
\end{itemize}

Теперь рассмотрим более сложные операции над кватернионами.Например умножение двух кватернионов. Такое умножение должно удовлетворять соотношениям Гамильтона:
\begin{equation*}
    \begin{array}{|c|c|c|c|}
        \hline
        \text{Произведение} & \text{Результат} & \text{Произведение} & \text{Результат} \\
        \hline
        i^2 & -\sigma & j^2 & -\sigma \\
        \hline
        k^2 & -\sigma & i j k & -\sigma \\
        \hline
        i j & k & j i & -k \\
        \hline
        j k & i & k j & -i \\
        \hline
        k i & j & i k & -j \\
        \hline
    \end{array}
\end{equation*}


Отметим также, что умножение не является коммутативным, то есть нам важен порядок умножения, но существуют условия коммутативности, описанные в ~\cite{Branets:1973}.Два кватерниона $p$ и $q$ коммутируют $(pq = qp)$ тогда и только тогда, когда выполняется хотя бы одно из условий:
\begin{enumerate}
    \item Если векторная часть одного из кватернионов равна 0, то есть или $p=0$, или $q=0$, такой кватернион является скаляром.
    \item У рассматриваемых кватернионов векторные части пропорциональны: $\vb{p} = \alpha \vb{q}$ для некоторого $\alpha \in \mathbb{R}$.
\end{enumerate}

Помимо умножения на действительное число, также есть общая формула умножения. Пусть $p = p_0\sigma + p_1 i + p_2 j + p_3 k$ и $q = q_0\sigma + q_1 i + q_2 j + q_3 k$, тогда с учетом таблицы для умножения базисных элементов получаем:
\begin{multline*}
    pq = (p_0\sigma + p_1 i + p_2 j + p_3 k)(q_0\sigma + q_1 i + q_2 j + q_3 k) = p_0 q_0 \sigma^2 + p_0 q_1 \sigma i + p_0 q_2 \sigma j + p_0 q_3 \sigma k + p_1 q_0 i \sigma + p_1 q_1 i^2 +\\
    + p_1 q_2 i j + p_1 q_3 i k + p_2 q_0 j \sigma + p_2 q_1 j i + p_2 q_2 j^2 + p_2 q_3 j k + p_3 q_0 k \sigma + p_3 q_1 k i + p_3 q_2 k j + p_3 q_3 k^2.
\end{multline*}
Далее сгруппируем по базисным элементам:
\begin{itemize}
    \item $p_0 q_0 - p_1 q_1 - p_2 q_2 - p_3 q_3$ --- коэффицент при $\sigma$;
    \item $p_0 q_1 + p_1 q_0 + p_2 q_3 - p_3 q_2$ --- коэффицент при $i$;
    \item $p_0 q_2 - p_1 q_3 + p_2 q_0 + p_3 q_1$ --- коэффицент при $j$;
    \item $p_0 q_3 + p_1 q_2 - p_2 q_1 + p_3 q_0$ --- коэффицент при $k$.
\end{itemize}

Есть более компактная формула умножения через скалярное и векторное умножения:
\begin{equation*}
    pq = (p_0 q_0 - (\vb{p}, \vb{q}))\sigma + p_0 \vb{q} + q_0 \vb{p} + \vb{p} \times \vb{q},
\end{equation*}
где $(\vb{p}, \vb{q}) = p_1 q_1 + p_2 q_2 + p_3 q_3$ --- скалярное произведение, а $\vb{p} \times \vb{q} = (p_2 q_3 - p_3 q_2)i + (p_3 q_1 - p_1 q_3)j + (p_1 q_2 - p_2 q_1)k$ --- векторное произведение.У скалярного умножения, также есть свойство инвариантности при циклической перестановке ~\cite{Branets:1973}:
\begin{equation*}
    \text{scal}(pqr) = \text{scal}(qrp) = \text{scal}(rpq),
\end{equation*}
где где $\text{scal}(q) = q_0$ --- скалярная часть кватерниона.

Также существует случай умножения двух чистых кватернионов.Пусть у нас есть два чистых кватерниона $p_0 = q_0 = 0$, тогда:
\begin{equation*}
    \vb{p} \vb{q} = -(\vb{p}, \vb{q})\sigma + \vb{p} \times \vb{q}.
\end{equation*}
Заметим также, что подобное умножение кватернионов связано с операциями векторной алгебры.Провести проверку можно с помощью базисных векторов. $ij = -1$, так как $(i,i)=1$ и $i \times i = 0$. Аналогично и для $i j = k$, так как $(i,j)=0$ и $i \times j = k$.

Умножение $pq = r$ можно записать в матричной форме. Если $r = r_0\sigma + r_1 i + r_2 j + r_3 k$, то:
\begin{equation*}
    \begin{bmatrix}
        r_0 \\ r_1 \\ r_2 \\ r_3
        \end{bmatrix}
        =
        \begin{bmatrix}
        p_0 & -p_1 & -p_2 & -p_3 \\
        p_1 &  p_0 & -p_3 &  p_2 \\
        p_2 &  p_3 &  p_0 & -p_1 \\
        p_3 & -p_2 &  p_1 &  p_0
        \end{bmatrix}
        \begin{bmatrix}
        q_0 \\ q_1 \\ q_2 \\ q_3.
        \end{bmatrix}
\end{equation*}

К основным свойствам умножения можно отнести:
\begin{itemize}
    \item Замкнутость --- в результате произведения двух кватернионов мы получим снова кватернион.
    \item Ассоциативность --- $(pq)r = p(qr)$, также проверяется прямым вычислением.
    \item Некоммутативность --- в общем случае $pq \neq qp$, происходит из-за векторного произведения.
    \item Дистрибутивность --- $p(q+r) = pq + pr$ и $(p+q)r = pr + qr$.
    \item Существование единицы --- это подразумевает, что есть кватернион($1 = \sigma + 0i + 0j + 0k$)  со скалярной частью равной единице и векторной частью равной нулю.
    \item Существование обратного --- это означает, что для любого ненулевого кватерниона существует $q^{-1}$ такой, что $q q^{-1} = q^{-1} q = \sigma$.
\end{itemize}

В алгебре кватернионов также вводится понятие как сопряжение. Для кватерниона $q = q_0 \sigma + \vb{q} = q_0 \sigma + q_1 i + q_2 j + q_3 k$ комплексно-сопряженный $q^*$ имеет вид:
\begin{equation*}
    q^* = q_0 \sigma - \vb{q} = q_0 \sigma - q_1 i - q_2 j - q_3 k.
\end{equation*}
У сопряжения также есть свой ряд свойств:
\begin{itemize}
    \item $(p + q)^* = p^* + q^*$ --- сопряжение суммы;
    \item $(pq)^* = q^* p^*$ --- сопряжение произведения, отметим, что поржок меняется;
    \item $q + q^* = 2q_0 \sigma$ --- сумма с сопряженным;
    \item $(q^*)^* = q$ --- дважды сопряженный.
\end{itemize}

Ранее уже была приведена формула нормы кватерниона ~\eqref{eq:(2)}, которая обозначается как $N(q)$ или $|q|$, рассмотрим вычисления подробнее:
\begin{multline*}
    |q|^2 = q^* q = (q_0 \sigma - \vb{q})(q_0 \sigma + \vb{q}) = q_0^2 \sigma^2 + q_0 \vb{q} - q_0 \vb{q} - \vb{q}\vb{q} = q_0^2 \sigma - \vb{q}\vb{q},
\end{multline*}
но в случае с чистым кватернионом будет верно $\vb{q}\vb{q} = -(\vb{q},\vb{q})\sigma = -|\vb{q}|^2 \sigma$, следовательно:
\begin{equation*}
    |q|^2 = q_0^2 \sigma + |\vb{q}|^2 \sigma = (q_0^2 + q_1^2 + q_2^2 + q_3^2)\sigma.
\end{equation*}
Таким образом мы приходим к формуле вида ~\eqref{eq:(2)}.

Отталкиваясь от понятия нормы кватерниона, можно перейти к тому какие кватернионы называют единичными или нормированными. Кватернион будет являться таковым, если $|q| = 1$, тогда нужно заметить, что все компоненты по модулю не превосходят 1.
Также существует норма произведения, обладающая свойством:
\begin{equation*}
    |pq| = |p| |q.
\end{equation*}
Для того, чтобы придти к следствию данного произведения необходимо доказательство:
\begin{equation*}
    |pq|^2 = (pq)(pq)^* = pq (q^* p^*) = p (q q^*) p^* = p |q|^2 p^* = p p^* |q|^2 = |p|^2 |q|^2.
\end{equation*}
Само следствие заключается в том, что в результате произведения двух единичных кватернионов получается снова единичный кватернион.

Обратный кватернион $q^{-1}$ определяется по формуле ~\eqref{eq:(3)}, но также существует условие существования такого кватерниона:
\begin{equation*}
    q q^{-1} = q^{-1} q = \sigma.
\end{equation*}
Сама же формула обратного кватерниона ~\eqref{eq:(3)} получается из определения нормы $q q^* = |q|^2 \sigma$.
Но существует и частный случай с единичным кватернионом. Если $|q| = 1$, то:
\begin{equation*}
    q^{-1} = q^*,
\end{equation*}
что аналогично ортогональным матрицам, где $A^{-1} = A^T$.

При работе с обратными кватернионами также важно помнить о порядке умножения так как они некоммутативны:
\begin{equation*}
    (pq)^{-1} = q^{-1} p^{-1},
\end{equation*}
потому что $(pq)(q^{-1}p^{-1}) = p(q q^{-1})p^{-1} = p p^{-1} = \sigma$.

\section{Пространственные перемещения с помощью кватернионов}

Ранее при описании операций над кватернионами уже рассматривались различные операции умножения с кватернионами, но не было до конца ясно, что именно происходит, когда один кватернион умножается на другой.Вращение в 4-мерном пространстве происходит в результате умножения на единичный кватернион $q = \cos\theta + \sin\theta\,\vb{n}$.
Для того, чтобы понять как именно это происходит необходимо разложить 4-мерное пространство на две взаимно перпендикулярные плоскости и посмотреть, что происходит в каждой из них.

Для любого единичного вектора $\vb{n}$ или если точнее чистого кватерниона с $\|\vb{n}\| = 1$ будет верно:
\begin{equation*}
    \vb{n}^2 = -\sigma.
\end{equation*}
Это происходит, потмоу что для чистого кватерниона выполняется: 
\begin{equation*}
    \vb{n}\vb{n} = -(\vb{n},\vb{n})\sigma + \vb{n}\times\vb{n} = -1\cdot\sigma + 0 = -\sigma.
\end{equation*}

В плоскости $\sigma, \vb{n}$ умножение на $(\vb{n}, \theta)$ будет аналогично умножению на $e^{i\theta}$ в комплексной плоскости, это является поворотом на угол $\theta$.Если $p$ лежит в плоскости $\sigma, \vb{n}$, то $q(\vb{n}, \theta) p$ поворот $p$ на $\theta$ и $p q(\vb{n}, \theta)$ также является поворотом на $\theta$, потому что в данной плоскости умножение коммутативно.

Любой вектор $\vb{n}$ перпендикулярный плоскости $\sigma, \vb{n}$, т.е. $(\vb{v}, \sigma) = 0$ и $(\vb{v}, \vb{n}) = 0$ будет верно, что $q(\vb{n}, \theta) \vb{v}$ остаенется в плоскости перпендикулярной $\sigma, \vb{n}$, но и в свою очередь $\vb{v} q(\vb{n}, \theta)$ также остается в этой плоскости. Это происходит, потому что умножение на единичный кватернион сохарняет углы, поэтому и ортогональность тоже сохранется.

Рассмотрим более подробно вращение в плоскости $\sigma, \vb{n}$. Пусть $\vb{n}$ является единичным вектором, $p$ кватернионов в плоскости $\sigma, \vb{n}$, т.е. $p = a\sigma + b\vb{n}$, тогда получатся следующие повороты:
\begin{enumerate}
    \item $L_{q}(p) = q p$ совершает поворот $p$ на угол $\theta$ в плоскости $\sigma, \vb{n}$.
    \item $R_{q}(p) = p q$ совершает поворот $p$ на угол $\theta$ в плоскости $\sigma, \vb{n}$.
    \item $R_{q^*}(p) = p q^*$ совершает поворот $p$ на угол $-\theta$ в плоскости $\sigma, \vb{n}$.
\end{enumerate}
Это доказывается тем, что $q^* = q(\vb{n}, -\theta)$ и свойством аддитивности угла:
\begin{equation*}
    q(\vb{n}, \theta) q(\vb{n}, \phi) = q(\vb{n}, \theta + \phi)
\end{equation*}

Если же мы хотим рассмотреть вращение в плоскости, перпендикулярной $\sigma, \vb{n}$, то нам необходим вектор $\vb{n}$, перпендикулярный нашей плоскости, т.е. мы берем чистый кватернион, который ортогонален $\vb{n}$, следовательно мы получим:
\begin{enumerate}
    \item $L_{q}(\vb{v}) = q \vb{v}$ совершает поворот $\vb{v}$ на угол $\theta$ в перпендикулярной плоскости.
    \item $R_{q}(\vb{v}) = \vb{v} q$ совершает поворот $\vb{v}$ на угол $-\theta$ в перпендикулярной плоскости.
    \item $R_{q^*}(\vb{v}) = \vb{v} q^*$ совершает поворот $\vb{v}$ на угол $\theta$ в перпендикулярной плоскости.
\end{enumerate}
Отметим, что направление вращения зависит от того, как именно мы умножаем, справа или слева.

Также в контексте пространственных перемещений существует ключевое тождество. Для любого комплексного числа $g = \gamma\sigma + \delta\vb{n}$ в плоскости $\sigma, \vb{n}$ будет верно:
\begin{equation*}
    \vb{v} g = g^* \vb{v}.
\end{equation*}

Для описания простых вращений используют сэндвич-операторы. Существуют основные используемые операторы такие, как оператор левого умножения $L_q(\rho) = q \rho$ и оператор правого умнлжения $R_q(\rho) = \rho q$. Для сэндвич-оператора с $q$ с двух сторон будет верно $T_q(\rho) = q \rho q$, а для оператора с сопряжением $S_q(\rho) = q \rho q^*$. Обозначения $L_q$ и $R_q$ --- это линейные преобразования(из-за дистрибутивности). В свою очередь $T_q$ и $S_q$ тоже являются линейными, но как композиции линейных.

Ранее мы рассматривали простую плоскость и перпендикулярную, сэндвич-операторы будут применимы и здесь. Для $p$ в плоскости $\sigma, \vb{n}$ произойдет следующее:
\begin{enumerate}
    \item $T_q(p) = q p q$ совершает поворот $p$ на угол $2\theta$ в плоскости $\sigma, \vb{n}$.
    \item $S_q(p) = q p q^*$ является тождественным преобразованием и поэтому оставляет $p$ неизменным.
\end{enumerate}

В перпендикулярной же плосокости $\sigma, \vb{n}$ для $\vb{v}$ преобразования получатся будут такими:
\begin{enumerate}
    \item $T_q(\vb{v}) = q \vb{v} q$ также является тождественным преобразованием.
    \item $S_q(\vb{v}) = q \vb{v} q^*$ совершает поворот $\vb{v}$ на угол $2\theta$.
\end{enumerate}
В результате, сэндвич-оператор совершает вращение так, что векторная часть остается всегда неизменной.

Существует также классификация подобных преобразований в пространостве. Аффиное сохраняет массу, обозначается как $S_q$,а полуаффинное сохраняет же нулевую массу и обозначается, как $T_N$. Также еще отмечают проективное, оно не сохраняет массу обозначается, как $T_{q(\vb{n},\theta)}$ при $\theta \neq \pi/2$. Разные операторы дают разные типы преобразований и нужно точно знать какой оператор использовать и для чего.

Начнем рассматривать различные преобразования с вращения. Сама задача может выглядеть, как то, что требуется повернуть некоторую точку $P$ или вектор $\vb{v}$ на угол $\theta$ вокруг оси $L$, которая к тому же будет проходить в направлении единичного вектора $\vb{n}$.
Любой вектор $\vb{v}$ можно разложить на две составляющие:
\begin{equation*}
    \vb{v} = \vb{v}_{\parallel} + \vb{v}_{\perp},
\end{equation*}
где $\vb{v}_{\parallel}$ является проекцией на ось вращения, а $\vb{v}_{\perp}$ перпендиклярная составляющая.
Само вращение будет происходить, как то, что паралелльная часть остается неизменной, а перпендикулярная совершает поворот в своей плоскости:
\begin{equation*}
    \vb{v}_{\perp} = \cos\theta\,\vb{v}_{\perp} + \sin\theta\,(\vb{n}\times\vb{v}_{\perp}).
\end{equation*}

Точка $P = O + (P - O)$ при приобразовании не теряет свое начало координат.

Вернемся к использованию сэндвич-операторов. Для единичного вектора $\vb{n}$ и вектора $\vb{v}$ используем оператор с сопряжением:
\begin{equation*}
    S_{q(\vb{n}, \theta/2)}(\vb{v}) = q(\vb{n}, \theta/2)\,\vb{v}\,q^*(\vb{n}, \theta/2),
\end{equation*}
здесь он поворачивает $\vb{v}$ на угол $\theta$ вокруг оси $\vb{n}$.

Аналогично можно использовать данный сэндвич-оператор и для поворота точек в пространстве. Для точки $P$  будет верно:
\begin{equation*}
    S_{q(\vb{n}, \theta/2)}(P) = q(\vb{n}, \theta/2)\,P\,q^*(\vb{n}, \theta/2),
\end{equation*}
здесь он поворачивает точку $P$ на угол $\theta$ вокруг оси $L$.

Существует также композиция вращений. Композиция двух вращений будет равна умножению соответствующих единичных кватернонов:
\begin{equation*}
    S_q \circ S_r = S_{qr}.
\end{equation*}
Множество вращений создает группу, изоморфную группе единичных кватернионов по умножению с точность до знака.

Если $q_1 = q(\vb{n}_1, \theta_1/2)$ и $q_2 = q(\vb{n}_2, \theta_2/2)$, то их произведение $q = q_1 q_2 = q(\vb{n}, \theta/2)$ дает вращение с параметрами:
\begin{itemize}
    \item Косинус половинного угла:
        \begin{equation*}
            \cos\frac{\theta}{2} = \cos\frac{\theta_1}{2}\cos\frac{\theta_2}{2} - \sin\frac{\theta_1}{2}\sin\frac{\theta_2}{2}\,(\vb{n}_1\cdot\vb{n}_2).
        \end{equation*}
    \item Ось вращения(умноженная на синус половинного угла):
        \begin{equation*}
            \sin\frac{\theta}{2}\,\vb{n} = \cos\frac{\theta_2}{2}\sin\frac{\theta_1}{2}\,\vb{n}_1 + \cos\frac{\theta_1}{2}\sin\frac{\theta_2}{2}\,\vb{n}_2 + \sin\frac{\theta_1}{2}\sin\frac{\theta_2}{2}\,(\vb{n}_1\times\vb{n}_2).
        \end{equation*}
    \item Случай, если оси совпадают $\vb{n}_1 = \vb{n}_2 = \vb{n}$, то углы просто складываются:
        \begin{equation*}
            q(\vb{n}, \theta_1/2) q(\vb{n}, \theta_2/2) = q(\vb{n}, (\theta_1+\theta_2)/2).
        \end{equation*}
\end{itemize}

Помимо вращения существует отражение объекта в пространстве. Например, требуется отразить точку $P$ или вектор $\vb{v}$ относительно плоскости $S$, которая также проходит через начало координат $O$ и перпендикулярна единичному вектору $\vb{n}$(нормали).

Как и для вращения можно разложить вектор $\vb{v} = \vb{v}_{\parallel} + \vb{v}_{\perp}$, здесь $\vb{v}_{\parallel}$ является частью паралелльную нормали(перпендикулярна плоскости), а $\vb{v}_{\perp}$ часть, лежащая в плоскости отражения.
После того, как отражение произошло, сама параллельная часть меняет знак на противположный. Перпендикулярная осатесят неизменной.

Для точки ничего не менется.

Существует, как и в случае с вращением, кватернионное представление. Оператор $T_{q(\vb{n},\theta)}$ является тождественным на векторах, которые перпендикулярны $\vb{n}$ и поворачивает векторы в плоскости $\sigma, \vb{n}$ на угол $2\theta$. Чтобы получить отражение, нужно, чтобы $\vb{v}_{\parallel}$ сменил свой знак, а $\vb{v}_{\perp}$ остался неизменным. Это будет соответствовать $2\theta = \pi$, т.е. $\theta = \pi/2$. Но $q(\vb{n}, \pi/2) = \cos(\pi/2)\sigma + \sin(\pi/2)\vb{n} = \vb{n}$.

Для того, чтобы отразить векторы в пространстве используем сэндвич-оператор с единичным вектором. Для единичного вектора $\vb{n}$ и вектора $\vb{vb}$ будет верно:
\begin{equation*}
    T_{\vb{n}}(\vb{v}) = \vb{n}\vb{v}\vb{n},
\end{equation*}
где оператор дает отражение $\vb{v}$ относительно плоскости, которая перпендикулярна $\vb{n}$.

В случае с отражениями также существует композиция. Композиция двух отражений эквивалентна одному вращению:
\begin{equation*}
    T_{\vb{n}_2} \circ T_{\vb{n}_1} = S_{\vb{n}_2\vb{n}_1},
\end{equation*}
но существует и обратное утверждение, что любое вращение может быть представлено как композиция двух отражений.

Если взять два единичных вектора $\vb{u}, \vb{v}$ в плоскости, которая перпендикулярна $\vb{n}$, таких что угол между ними равен $\theta/2$, то;
\begin{equation*}
    -\vb{v}\vb{u} = \cos(\theta/2)\sigma + \sin(\theta/2)\vb{n} = q(\vb{n}, \theta/2),
\end{equation*}
\begin{equation*}
    S_{q(\vb{n}, \theta/2)} = T_{\vb{v}} \circ T_{\vb{u}}.
\end{equation*}

Заметим, что сэндвич-оператор с $\vb{n}$ не дает отражение точки. Для точки $P$:
\begin{equation*}
    T_{\vb{n}}(P) = \vb{n} P \vb{n},
\end{equation*}
где сэндвич-оператор поворачивает точку на угол $\pi$ вокруг оси $\vb{n}$, а не отражается относительно плоскостию Но возникает тогда вопрос, как же правильно получить отражение точки? Чтобы получить отражение точки $P$ относительно плоскости через $O$ перпендикулярно $\vb{n}$, нужно применять сэндвич-оператор к $P - 2\sigma$:
\begin{equation*}
    T_{\vb{n}}(P - 2\sigma) = \vb{n}(P - 2\sigma)\vb{n}.
\end{equation*}

Еще не менее важная тема в контексте преобразований это проекция. Например, мы хотим спроецировать точку $P$ из некоторой точки обзора $E$ на плоскость проекции $S$, которая проходит через точку перпендикулярную единичному вектору нормали $\vb{n}$.В данном случае мы можем использовать оператор $T_{q(\vb{n}, -\pi/4)}$. Использование $\theta = -\pi/4$ является наиболее подходящим, потому что $T_{q(\vb{n}, -\pi/4)}$ вращает плоскость $\sigma, \vb{n}$ на угол $2(-\pi/4) = -\pi/2$, а это значит, что единичный вектор $\vb{n}$ переходит в точку $\sigma$, т.е. начало координат. Начало координа же в свою очередь переходит в вектор $\vb{n}$. Таким образом, оператор как бы преобразует расстояние вдоль нормали в массу в начале координат.

Рассмотрим ортогональные матрицы и их связь с перспективной проекцией. Предположим, что точка проекции на оси $Oz$ будет $E = (0, 0, 1)$, плоскость проекции будет $XY$- плоскостью, т.е. $z=0$, а нормаль $\vb{n} = (0, 0, -1)$.
Для точки $P = (x,y,z)$ формула перспективной проекции будет выглядеть следующим образом:
\begin{equation*}
    G = \left( \frac{x}{1-z}, \frac{y}{1-z}, 0 \right).
\end{equation*}
Сама матрица проекции $4\times 4$:
\begin{equation*}
    M = \begin{pmatrix}
        1 & 0 & 0 & 0 \\
        0 & 1 & 0 & 0 \\
        0 & 0 & 0 & -1 \\
        0 & 0 & 0 & 1
        \end{pmatrix}
\end{equation*}
Применим теперь кватернионный подход, сэндвич-оператор $T_{q(\vb{n}, -\pi/4)}$ является вращением в 4-мерном пространстве в плоскости $\sigma, \vb{n}$. В нашем случае плоскость $\sigma, \vb{n}$ --- это плоскость $w(-z)$, т.е. вращение на $-\pi/2$ в плоскости $zw$.Тогда соответствующая ортогональная матрица:
\begin{equation*}
    R = \begin{pmatrix}
        1 & 0 & 0 & 0 \\
        0 & 1 & 0 & 0 \\
        0 & 0 & 0 & -1 \\
        0 & 0 & 1 & 0
        \end{pmatrix}.
\end{equation*}
Перспективную проекцию можно реализовать как вращение в 4-мерном пространстве с последующим ортогональным проецированием.

В контексте применения к точкам, $T_{q(\vb{n}, -\pi/4)}(P)$ не дает перспективную проекци точки. Оператор не является аффинным, он переводит начало координат в вектор:
\begin{equation*}
    T_{q(\vb{n}, -\pi/4)}(\sigma) = -\vb{n} \neq \sigma,
\end{equation*}
но $T_{q(\vb{n}, -\pi/4)}(P) - T_{q(\vb{n}, -\pi/4)}(E) = T_{q(\vb{n}, -\pi/4)}(P - E)$, поэтому с помощью вычитания точечной массы, мы все еще можем получить проекцию.

Также отметим, что перенос и проекция коммутируют. Проецирование точки из точки проекции $E$ на плоскость $S$ с последующим переносом на вектор $\vb{v}$ является эквивалентным сначала переносу всего на $\vb{v}$, а затем проецированию из перенесенной точки проекции $E + \vb{v}$ на перенесенную плоскость $S + \vb{v}$.

Ранее были изучены уже два сэндвич-оператора $S_q(\rho) = q\rho q^*$, которое является вращением в 3-мерном пространстве и $T_q(\rho) = q\rho q$ отражение(при $q = \vb{n}$) или перспективная проекция(при $q = q(\vb{n}, -\theta/2)$).Но мы не рассматривалии более подробно простые умножения вроде левого $L_q(\rho) = q\rho$ и правого $R_q(\rho) = \rho q$. Они представляют собой двойные вращения в 4-мерном пространстве и вращают одновременно в обеих плоскостях. В 3-мерном пространстве это будет давать комбинацию вращения и проекции.

Пусть $r = q(\vb{n}, \theta/2)$, тогда $r^2 = q$ и $r^* = r^{-1}$. Отметим также, что в источнике ~\cite{Goldman:2010} приведено разложение вида $q\rho = r(r\rho r^*) = T_r(S_r(\rho))$, следовательно:
\begin{equation*}
    L_q(\rho) = T_r(S_r(\rho)) = S_r(T_r(\rho)).
\end{equation*}
В общем случае они не коммутируют, но в данном да,  потому что  $r$ коммутирует с $S_r(\rho)$ в плоскости $\sigma, \vb{n}$, тогда можно сделать вывод, что левое умножение на $q$ явялется комбинацией вращения $S_r$ и перспективной проекции $T_r$. Такие преобразования и называются роторперспективами.

Есть различные частные случаи, например, когда $\theta = \pi/2$. Тогда $q = \vb{n}$ и $r = q(\vb{n}, \pi/4)$. Сначала нужно повернуть вектор $\vb{v}$ на угол $\pi/2$ вокруг оси $\vb{n}$"
\begin{equation*}
    S_r(\vb{v}) = (\vb{v}\cdot\vb{n})\vb{n} + \vb{n}\times\vb{v},
\end{equation*}
а затем спроецировать полученную точку $\sigma + S_r(\vb{v})$ из точки проекции $E = \sigma + \vb{n}$ на плоскость через $\sigma$, перпендикулярную $\vb{n}$:
\begin{equation*}
    T_r(S_r(\vb{v})) = -(\vb{n}\cdot\vb{v})\sigma + \vb{n}\times\vb{v} = \vb{n}\vb{v}.
\end{equation*}
Таким образом, произведение двух векторов $\vb{n}\vb{v}$ и есть роторперспектива.

Есть случай и при $\theta = \pi$, $q = -\sigma$, $r = q(\vb{n}, \pi/2) = \vb{n}$, тогда $S_r$ вляется вращением на $\pi$ вокруг $\vb{n}$, а $T_r = T_{\vb{n}}$ отражением в плоскости, перпендикулярной $\vb{n}$. Комбинация вращения и отражения называется ротоотражением.

